%% ============================================================
%% BÖLÜM 17 — Parametresiz İstatistik  [OPSİYONEL]
%% Olasılık Teorisi ve Matematiksel İstatistik
%% ============================================================
\chapter{Parametresiz İstatistik}
\label{chp:parametresiz}

\bolumalinti{%
  The great advantage of the nonparametric approach is that we need not make
  any assumption about the form of the underlying population.
}{Frank Wilcoxon}

\begin{kazanimlar}
Bu bölümün sonunda okuyucu:
\begin{itemize}
  \item Parametresiz yöntemlerin motivasyonunu ve güç-sağlamlık dengesini açıklar.
  \item İşaret testi ve Wilcoxon işaret-sıra testini uygular ve asimptotik dağılımlarını türetir.
  \item Mann-Whitney $U$-testini iki bağımsız örneklem karşılaştırması için kullanır.
  \item Kruskal-Wallis ve Friedman testlerini çok grup karşılaştırmasında uygular.
  \item Kolmogorov-Smirnov testini dağılım uyum testi olarak kullanır.
  \item Bootstrap yöntemini standart hata ve güven aralığı hesabında uygular.
\end{itemize}
\end{kazanimlar}

\begin{motivasyon}
Normal dağılım varsayımı zaman zaman geçersizdir: çarpık dağılımlar, aykırı değerler, sıralı ölçek. Parametresiz yöntemler dağılım hakkında minimal varsayımla çalışır. "Sağlamlık" (robustness) modern istatistikte kritik önem taşır.

\medskip
\noindent\textbf{Köprü:} Bölüm~\ref{chp:hipotez_testi}'deki test çerçevesi burada da geçerlidir; fark, test istatistiklerinin sıralardan ya da işaretlerden inşa edilmesidir.
\end{motivasyon}

%% ============================================================
\section{İşaret Testi}
\label{sec:isaret_testi}
%% ============================================================

\begin{kesif}
Sürekli dağılım medyanı $m$ için $X_i>m$ veya $X_i<m$ ikili bilgisi kullanılabilir. $H_0:m=m_0$ altında bu ikili değerlerin dağılımı nedir?
\end{kesif}

\begin{tanim}[İşaret testi]\label{def:isaret_testi}
$X_1,\ldots,X_n\iid F$ sürekli. $H_0: F(m_0)=1/2$ (medyan $m_0$). $B = \#\{i:X_i>m_0\}$. $H_0$ altında $B\sim\Binom(n,1/2)$.

Test istatistiği $B$ (veya $\min(B,n-B)$ iki yönlü test için). Küçük $n$'de kesin Binom dağılımı; büyük $n$'de $B$ yaklaşık $\Normal(n/2,n/4)$.
\end{tanim}

\begin{teorem}[İşaret testinin asimptotik dağılımı]
$H_0$ altında $Z=(2B-n)/\sqrt{n}\dto\Normal(0,1)$.
\end{teorem}

\begin{proof}
$B=\sum_{i=1}^n\mathbf{1}[X_i>m_0]$; her terim i.i.d.\ $\mathrm{Ber}(1/2)$. MLT uygulanır.
\end{proof}

%% ============================================================
\section{Wilcoxon İşaret-Sıra Testi}
\label{sec:wilcoxon_isaret}
%% ============================================================

\begin{tanim}[Wilcoxon işaret-sıra istatistiği]\label{def:wilcoxon_isaret}
$X_1,\ldots,X_n$ simetrik sürekli dağılımdan (simetri merkezi $\theta$), $H_0:\theta=0$. $|X_1|,\ldots,|X_n|$'yi küçükten büyüğe sırala; $R_i = |X_i|$'nin sırası.

\textbf{Wilcoxon istatistiği:} $W^+ = \sum_{i=1}^n R_i\,\mathbf{1}[X_i>0]$.
\end{tanim}

\begin{teorem}[Wilcoxon asimptotik dağılımı]\label{thm:wilcoxon_asimptotik}
$H_0$ altında:
\[
  \Beklenti[W^+] = \frac{n(n+1)}{4}, \qquad \Var(W^+) = \frac{n(n+1)(2n+1)}{24}.
\]
Standartlaştırılmış $Z = (W^+ - n(n+1)/4)/\sqrt{n(n+1)(2n+1)/24} \dto \Normal(0,1)$.
\end{teorem}

\begin{proof}[Proof — beklenti]
$H_0$ altında her $X_i$ için $\Prob(X_i>0)=1/2$ (simetri). $W^+=\sum_{i=1}^n R_i\,\psi_i$ ($\psi_i=\mathbf{1}[X_i>0]$). $\psi_i\perp R_i$ (simetri) ve $\Beklenti[\psi_i]=1/2$.
\[
  \Beklenti[W^+] = \sum_{i=1}^n\Beklenti[R_i]\cdot\Beklenti[\psi_i] = \frac{1}{2}\sum_{i=1}^n i = \frac{n(n+1)}{4}.
\]
\end{proof}

\begin{ispatiskele}
Varyans hesabı: $\psi_i,\psi_j$ ve $R_i,R_j$ bağımsız olduğundan $\Cov(R_i\psi_i,R_j\psi_j)=0$ ($i\neq j$). $\Var(R_i\psi_i)=\Beklenti[R_i^2]/4$; $\sum_{i=1}^n i^2=n(n+1)(2n+1)/6$. Bölünce $n(n+1)(2n+1)/24$.
\end{ispatiskele}

%% ============================================================
\section{Mann-Whitney $U$-Testi}
\label{sec:mann_whitney}
%% ============================================================

\begin{tanim}[Mann-Whitney istatistiği]\label{def:mann_whitney}
$X_1,\ldots,X_m\iid F$ ve $Y_1,\ldots,Y_n\iid G$ bağımsız. $H_0:F=G$.

$U = \#\{(i,j): X_i < Y_j\} = \sum_{i=1}^m\sum_{j=1}^n\mathbf{1}[X_i<Y_j]$.
\end{tanim}

\begin{teorem}[Mann-Whitney asimptotik dağılımı]\label{thm:mann_whitney}
$H_0$ altında:
\[
  \Beklenti[U] = \frac{mn}{2}, \qquad \Var(U) = \frac{mn(m+n+1)}{12}.
\]
$Z = (U - mn/2)/\sqrt{mn(m+n+1)/12}\dto\Normal(0,1)$.
\end{teorem}

\begin{proof}[Proof — beklenti]
$H_0:F=G$ altında $\Prob(X_i<Y_j)=1/2$ (sürekli dağılım). $U=\sum_{ij}\mathbf{1}[X_i<Y_j]$, her biri $\mathrm{Ber}(1/2)$; $\Beklenti[U]=mn/2$.
\end{proof}

\begin{teorem}[Wilcoxon sıra-toplam — eşdeğerlik]\label{thm:wilcoxon_mann}
İki örneklem birleştirilip sıralanır; $W_X = $ X değerlerinin sıra toplamı. $U = W_X - m(m+1)/2$ (sıra toplamından Mann-Whitney elde edilir). İki test eşdeğerdir.
\end{teorem}

%% ============================================================
\section{Kruskal-Wallis ve Friedman Testleri}
\label{sec:kruskal_friedman}
%% ============================================================

\begin{tanim}[Kruskal-Wallis testi]\label{def:kruskal_wallis}
$k$ grup, grup $i$'de $n_i$ gözlem; toplam $N=\sum n_i$. Tüm gözlemler birleştirilip sıralanır; $\bar{R}_i$ = $i$. grubun ortalama sırası. Kruskal-Wallis istatistiği:
\[
  H = \frac{12}{N(N+1)}\sum_{i=1}^k n_i(\bar{R}_i - \bar{R})^2
\]
$\bar{R}=(N+1)/2$ (genel ortalama sıra).
\end{tanim}

\begin{teorem}[Kruskal-Wallis asimptotiği]\label{thm:kruskal_asimptotik}
$H_0$ (tüm gruplar aynı dağılım) altında $H\dto\chi^2_{k-1}$.
\end{teorem}

\begin{ispatiskele}
ANOVA sıra dönüşümü: sıraların $F$-istatistiği formülü uygulanınca $H$ elde edilir; $\chi^2$ asimptotiği normal yakınsama argümanından gelir.
\end{ispatiskele}

\begin{tanim}[Friedman testi]\label{def:friedman}
$n$ blok, $k$ muamele; blok içi sıralar $r_{ij}$ ($j=1,\ldots,k$). Friedman istatistiği:
\[
  F_r = \frac{12n}{k(k+1)}\sum_{j=1}^k\!\left(\bar{r}_j - \frac{k+1}{2}\right)^{\!2}
\]
$\bar{r}_j$ = $j$-inci muamelenin ortalama sırası. $H_0$ altında $F_r\dto\chi^2_{k-1}$.
\end{tanim}

%% ============================================================
\section{Kolmogorov-Smirnov Testi}
\label{sec:ks_testi}
%% ============================================================

\begin{tanim}[Ampirik dağılım fonksiyonu]\label{def:ampirik_df}
$X_1,\ldots,X_n$ gözlemi için \textbf{ampirik dağılım fonksiyonu} \emph{(empirical CDF)}:
\[
  F_n(x) = \frac{1}{n}\sum_{i=1}^n\mathbf{1}[X_i\leq x].
\]
\end{tanim}

\begin{teorem}[Glivenko-Cantelli]\label{thm:glivenko_cantelli}
$X_i\iid F$. O zaman
\[
  \sup_x |F_n(x)-F(x)| \asto 0.
\]
\end{teorem}

\begin{proof}[Proof (taslak)]
Her sabit $x$ için GBSK'dan $F_n(x)\asto F(x)$. Üst limit geçirgenliği (üstel sınır argümanı) sayılabilir yoğun $\{x_k\}$ kümesi üzerinden $\sup$'u kapsar; sol süreklilik ve sınırlılık ile genel sonuç elde edilir.
\end{proof}

\begin{tanim}[KS istatistiği]\label{def:ks_istatistik}
$H_0:F=F_0$ (belirtilmiş sürekli dağılım). \textbf{Kolmogorov-Smirnov istatistiği:}
\[
  D_n = \sup_x |F_n(x) - F_0(x)|.
\]
\end{tanim}

\begin{teorem}[Kolmogorov dağılımı]\label{thm:kolmogorov_dagilim}
$H_0$ altında $F_0$ sürekli ise:
\[
  \sqrt{n}\,D_n \dto K
\]
$K$'nın dağılım fonksiyonu: $\Prob(K\leq t) = \sum_{k=-\infty}^{\infty}(-1)^k e^{-2k^2t^2}$.
\end{teorem}

\begin{ispatiskele}
$U_i = F_0(X_i)\iid\Uniform[0,1]$ ($F_0$ sürekli). $F_n$ yerine uniform ampirik süreç $\sqrt{n}(\hat{F}_n(u)-u)$ incelenir; bu Brownian köprüye yakınsar (Donsker). Brownian köprünün sonsuz normu Kolmogorov dağılımını verir.
\end{ispatiskele}

%% ============================================================
\section{Medyan için Güven Aralığı}
\label{sec:medyan_ga}
%% ============================================================

\begin{teorem}[Medyan için dağılım-serbest GA]\label{thm:medyan_ga}
$X_1,\ldots,X_n\iid F$ sürekli; medyan $m$ bilinmiyor. Sıra istatistikleri $X_{(1)}\leq\cdots\leq X_{(n)}$.

$\Prob(X_{(r)}\leq m\leq X_{(s)}) = \sum_{k=r}^{s-1}\binom{n}{k}(1/2)^n$.

$(1-\alpha)$ güven aralığı: $r$ ve $s$'yi $\sum_{k=r}^{s-1}\binom{n}{k}2^{-n}\geq1-\alpha$ sağlayacak biçimde seç.
\end{teorem}

\begin{proof}
$\Prob(X_{(r)}\leq m) = \Prob(\text{en az } r \text{ gözlem } m\text{'den küçük}) = \Prob(B_{n,1/2}\geq r)$ ($B_{n,1/2}\sim\Binom(n,1/2)$). $\Prob(X_{(r)}\leq m\leq X_{(s)}) = \Prob(r\leq B_{n,1/2}\leq s-1)$.
\end{proof}

\begin{ornek}[Medyan GA, $n=9$, $\alpha=0.05$]
$\sum_{k=r}^{8-r}\binom{9}{k}(1/2)^9$ $r=2$, $s=8$: $\sum_{k=2}^7\binom{9}{k}/512=(9+36+84+126+126+84)/512=465/512=0.908<0.95$. $r=2,s=8$ yeterli değil. $r=1,s=9$: $\sum_{k=1}^8\binom{9}{k}/512=(512-1-1)/512=510/512=0.996\geq0.95$. GA: $[X_{(1)},X_{(9)}]$, kapsama $\geq99.6\%$.
\end{ornek}

%% ============================================================
\section{Spearman Sıra Korelasyonu}
\label{sec:spearman}
%% ============================================================

\begin{tanim}[Spearman'ın $\rho$ istatistiği]\label{def:spearman}
$(X_i,Y_i)_{i=1}^n$ çiftleri. $R_i$ = $X_i$'nin sırası, $S_i$ = $Y_i$'nin sırası. \textbf{Spearman $\rho$}:
\[
  r_S = 1 - \frac{6\sum_{i=1}^n d_i^2}{n(n^2-1)}, \quad d_i = R_i - S_i.
\]
Parametrik Pearson korelasyonunun sıra tabanlı alternatifi; aykırı değerlere sağlamlı.
\end{tanim}

\begin{teorem}[Spearman $\rho$ — bağlantı]
$r_S$ aslında sıraların Pearson korelasyonudur: $r_S = \Cor(R_i,S_i)$. $H_0:\rho=0$ (bağımsızlık) altında $T=r_S\sqrt{(n-2)/(1-r_S^2)}\dto t_{n-2}$.
\end{teorem}

\begin{proof}
Sıralar $1,\ldots,n$'nin ortalaması $(n+1)/2$; varyansı $(n^2-1)/12$. Pearson formülünü sıralara uygulayınca $r_S$ elde edilir. Büyük örneklem: $\sqrt{n-1}r_S\dto\Normal(0,1)$ ($H_0$ altında).
\end{proof}

\begin{ornek}[Spearman $\rho$ hesabı]
$(X,Y)$ çiftleri: sıralar $(R,S)$: $(1,2),(2,1),(3,4),(4,3),(5,5)$. $d_i=(-1,1,-1,1,0)$. $\sum d_i^2=4$. $r_S=1-6\cdot4/(5\cdot24)=1-0.2=0.8$.
\end{ornek}

%% ============================================================
\section{Permütasyon Testleri}
\label{sec:permutasyon_test}
%% ============================================================

\begin{tanim}[Permütasyon testi]\label{def:permutasyon_test}
$H_0$: iki grup $\{X_i\}$ ve $\{Y_j\}$ aynı dağılımdan. Test istatistiği $T$ (örn. $|\bar{X}-\bar{Y}|$).

\textbf{Permütasyon dağılımı:} Tüm $\binom{m+n}{m}$ mümkün grup atamasında $T$ değerleri hesaplanır. $p$-değeri: gözlenen $T$'den büyük ya da eşit $T$ değerlerinin oranı.
\end{tanim}

\begin{teorem}[Permütasyon testinin kesin geçerliliği]
$H_0$ altında tüm permütasyonlar eşit olasılıklı; permütasyon $p$-değeri tam anlamlılık düzeyi $\alpha$'yı kontrol eder (dağılım varsayımı gerektirmez).
\end{teorem}

\begin{ornek}[Küçük örneklem permütasyon testi]
Grup A: $\{3,7\}$, Grup B: $\{1,5\}$. $T=|\bar{A}-\bar{B}|=|5-3|=2$. Tüm $\binom{4}{2}=6$ bölme: $\{3,7\}$ vs $\{1,5\}$: $T=2$; $\{3,1\}$ vs $\{7,5\}$: $T=4$; $\{3,5\}$ vs $\{1,7\}$: $T=2$; $\{1,7\}$ vs $\{3,5\}$: $T=2$; $\{1,5\}$ vs $\{3,7\}$: $T=2$; $\{7,5\}$ vs $\{3,1\}$: $T=4$. $p=\Prob(T\geq2)=6/6=1.00$. $H_0$ reddedilemez.
\end{ornek}

%% ============================================================
\section{Asimptotik Göreli Verimlilik (ARE)}
\label{sec:are}
%% ============================================================

\begin{tanim}[ARE]\label{def:are}
$T_n^{(1)}$ ve $T_n^{(2)}$ aynı hipotez için iki test. $n_1(n_2)$: $T^{(1)}(T^{(2)})$'nin belirli güce ($1-\beta$) ulaşması için gereken örneklem büyüklükleri.
\[
  \mathrm{ARE}(T^{(1)},T^{(2)}) = \lim_{n\to\infty}\frac{n_2}{n_1}.
\]
$\mathrm{ARE}>1$: $T^{(1)}$ daha verimli; daha az gözlemle aynı güce ulaşır.
\end{tanim}

\begin{teorem}[Wilcoxon vs $t$-testi ARE]\label{thm:are_wilcoxon}
Normal dağılım altında:
\[
  \mathrm{ARE}(\text{Wilcoxon},t) = \frac{3}{\pi} \approx 0.955.
\]
Wilcoxon normallik geçerliyken $t$'nin $\%95.5$ veriminde çalışır — ama normallikten sapmalarda Wilcoxon üstün olabilir.
\end{teorem}

\begin{ispatiskele}
Pitman verimlilik: $\mathrm{ARE}=\mathrm{SH}^2_{t}/\mathrm{SH}^2_W$ lokal alternatif altında asimptotik güç eğrileri karşılaştırılarak türetilir. Normal $f$ için Pitman verimlilik $12\sigma^2[f(0)]^2 = 12\sigma^2 /(2\pi\sigma^2)=6/\pi\to3/\pi$ (tek örneklem Wilcoxon için $3/\pi\approx0.955$).
\end{ispatiskele}

\begin{teorem}[ARE — genel sonuç]
\begin{itemize}
  \item Herhangi sürekli dağılım için: $\mathrm{ARE}(\text{Wilcoxon},t)\geq108/125\approx0.864$.
  \item Çift üstel dağılımda: $\mathrm{ARE}=2.0$ (Wilcoxon iki kat verimli).
  \item Cauchy dağılımda: $\mathrm{ARE}=\infty$ ($t$-testi tutarsız; Wilcoxon tutarlı).
\end{itemize}
Sonuç: Wilcoxon normalliği hiçbir zaman çok fazla kaybetmez; ağır kuyrukta büyük kazanç.
\end{teorem}

%% ============================================================
\section{Bootstrap}
\label{sec:bootstrap}
%% ============================================================

\begin{kesif}
Tahmin edicimizin standart hatasını analitik olarak hesaplayamıyorsak ne yaparız? Bootstrap fikri: gerçek popülasyondan örnekleme yerine, elimizdeki veriyi popülasyon gibi kullan ve yeniden örnekle.
\end{kesif}

\begin{tanim}[Bootstrap]\label{def:bootstrap}
$\mathbf{x}=(x_1,\ldots,x_n)$ gözlem. Bootstrap algoritması:
\begin{enumerate}
  \item $b=1,\ldots,B$ için: $\mathbf{x}^{*b}=(x_1^*,\ldots,x_n^*)$ yerine koyarak $\mathbf{x}$'den çek.
  \item $\hat\theta^{*b} = T(\mathbf{x}^{*b})$ hesapla.
  \item Bootstrap standart hata: $\hat{\mathrm{SE}}_B = \sqrt{\frac{1}{B-1}\sum_{b=1}^B(\hat\theta^{*b}-\bar{\hat\theta}^*)^2}$.
\end{enumerate}
\end{tanim}

\begin{tanim}[Bootstrap güven aralıkları]\label{def:bootstrap_ga}
\begin{itemize}
  \item \textbf{Yüzdelik GA:} $[\hat\theta^*_{(\alpha/2)}, \hat\theta^*_{(1-\alpha/2)}]$.
  \item \textbf{BCa (bias-corrected and accelerated):} Yüzdelik yöntemini çarpıklık ve ivme için düzeltir.
\end{itemize}
\end{tanim}

\begin{teorem}[Bootstrap tutarlılığı]\label{thm:bootstrap_tutarlilik}
Geniş koşullar altında bootstrap dağılımı, gerçek örnekleme dağılımına yakınsar. Özellikle $\sqrt{n}$-tutarlı ve asimptotik normal $T(\mathbf{X})$ için bootstrap standart hata tutarlıdır.
\end{teorem}

\begin{ispatiskele}
Genelleşmiş delta yöntemi: $T(\mathbf{X})=g(\bar{X})$ biçimindeyse, $T^* - T$ dağılımının, $T - \theta$ dağılımına bootstrap yakınsaması MLT + Slütseva argümanıyla kurulur.
\end{ispatiskele}

\begin{disiplinlerarasibaglan}
\textbf{Biyoistatistik:} Küçük klinik örneklemlerde bootstrap güven aralıkları standart. \textbf{Makine öğrenmesi:} Bootstrap + bagging ile varyans azaltma (rastgele ormanlar). \textbf{Ekonometri:} Küme-sağlamlı standart hatalar bootstrap ile hesaplanır.
\end{disiplinlerarasibaglan}

%% ============================================================

\begin{mikroozet}
\begin{itemize}
  \item Parametresiz testler: minimal dağılım varsayımı; sıra veya işaret tabanlı.
  \item İşaret: medyan testi; Wilcoxon: simetri testi; Mann-Whitney: iki örneklem; KW/Friedman: çok grup.
  \item KS testi: ampirik ile teorik KDF farkı; $\sqrt{n}D_n\to$ Kolmogorov.
  \item Bootstrap: yeniden örnekleme ile standart hata ve GA; dağılım varsayımı gerektirmez.
\end{itemize}
\end{mikroozet}

\begin{ozyansima}
Parametresiz testler normallik varsayımı bozulduğunda her zaman parametrik testlerden daha iyi midir? Normallik geçerliyken ne kadar güç kaybı olur?
\end{ozyansima}

\begin{kopru}
\textbf{Nereden geldik:} Bölüm~\ref{chp:hipotez_testi}'deki test çerçevesi parametresiz yöntemleri de kapsar; temel fark test istatistiğinin seçimindedir.

\textbf{Nereye gidiyoruz:} Bölüm~\ref{chp:asimptotik}'da bootstrap tutarlılığının arkasındaki asimptotik teoriye, fonksiyonel MLT'ye ve Donsker teoreminin kesin formülasyonuna bakacağız.
\end{kopru}

%% ============================================================
%% ALISTIRMALAR
%% ============================================================

\cozumlualistirmalar

\begin{alistirma}[Al.17.1]{A}{İşaret testi}
$n=15$ hastada plasebo ile karşılaştırılan ilacın etkisi: 11 hasta iyileşti, 4 kötüleşti. $H_0:m=0$ (medyan fark sıfır) vs $H_1:m>0$. $\alpha=0.05$ için karar ver.
\end{alistirma}
\alcozum{%
$B=11\sim\Binom(15,0.5)$ ($H_0$ altında). $p=\Prob(B\geq11)=\sum_{k=11}^{15}\binom{15}{k}(0.5)^{15}=(1+15+105+455+3003\cdot0)/32768$. Doğrudan: $\Prob(B\geq11)=0.059>0.05$. $H_0$ reddedilemez (tam Binom ile $p=0.059$).%
}

\begin{alistirma}[Al.17.2]{B}{Mann-Whitney}
Grup A: $\{3,5,7,9,11\}$, Grup B: $\{2,4,8,10,12\}$. $H_0:F_A=F_B$. $U$ istatistiğini hesapla ve asimptotik $z$ ile test yap.
\end{alistirma}
\alcozum{%
$U=\#\{(A_i,B_j):A_i<B_j\}$. $3<4,8,10,12$ (4), $5<8,10,12$ (3), $7<8,10,12$ (3), $9<10,12$ (2), $11<12$ (1): $U=13$. $\Beklenti=12.5$, $\Var=5\cdot5\cdot11/12=22.9$. $Z=(13-12.5)/\sqrt{22.9}=0.1$. $H_0$ reddedilemez.%
}

\alistirmalar

\begin{alistirma}[Al.17.3]{A}{KS testi uygulaması}
$n=5$ gözlem: $\{0.2, 0.5, 0.6, 0.8, 0.9\}$. $H_0: F_0 = \Uniform[0,1]$. $D_5$'i hesapla.
\end{alistirma}
\alcozum{%
$F_5(x_i)=i/5$, $F_0(x_i)=x_i$. Farklar: $|0.2-0.2|=0$, $|0.4-0.5|=0.1$, $|0.6-0.6|=0$, $|0.8-0.8|=0$, $|1.0-0.9|=0.1$. Sol sınır farkları da hesapla: max$=0.1$. $D_5=0.1$. $\sqrt{5}\cdot0.1=0.224$ — KS tablosundan $\alpha=0.05$ kritik değeri $0.565$; $H_0$ reddedilemez.%
}

\begin{alistirma}[Al.17.4]{B}{Bootstrap standart hata}
$n=6$ gözlem: $\{2,4,5,7,9,12\}$. Medyanın bootstrap standart hatasını $B=5$ bootstrap yinelemesiyle elle hesapla (rastgele çekimler verildiğinde: $\{2,4,4,7,9,9\}$, $\{4,5,7,7,12,12\}$, $\{2,2,5,9,9,12\}$, $\{4,5,7,9,9,12\}$, $\{2,4,5,5,7,12\}$).
\end{alistirma}
\alcozum{%
Medyan(orijinal)$=6$. Bootstrap medyanları: $m_1=(4+7)/2=5.5$, $m_2=(7+7)/2=7$, $m_3=(5+9)/2=7$, $m_4=(7+9)/2=8$, $m_5=(5+5)/2=5$. $\bar{m}^*=6.5$. $\hat{SE}=\sqrt{[(5.5-6.5)^2+(7-6.5)^2+(7-6.5)^2+(8-6.5)^2+(5-6.5)^2]/4}=\sqrt{(1+0.25+0.25+2.25+2.25)/4}=\sqrt{1.5}=1.22$.%
}

\begin{alistirma}[Al.17.5]{C}{Glivenko-Cantelli hızı}
$X_1,\ldots,X_n\iid\Uniform[0,1]$. $\Beklenti[\sup_x|F_n(x)-x|]$'in üst sınırını Dvoretzky-Kiefer-Wolfowitz eşitsizliğini kullanarak ver: $\Prob(\sup_x|F_n(x)-F(x)|>\varepsilon)\leq 2e^{-2n\varepsilon^2}$.
\end{alistirma}
\alcozum{%
DKW'den $\Prob(D_n>\varepsilon)\leq2e^{-2n\varepsilon^2}$. $\Beklenti[D_n]=\int_0^\infty\Prob(D_n>t)\,dt\leq\int_0^\infty2e^{-2nt^2}\,dt=2\cdot\frac{\sqrt{\pi}}{2\sqrt{2n}}=\sqrt{\pi/(2n)}$. Dolayısıyla $\Beklenti[D_n]=O(1/\sqrt{n})$.%
}
